%=============================================================================
% NEW PATENT PROPOSALS
% Identified from Patent #11 EM Coupling SRF research
% Three new applications covering gaps in the existing DFD patent portfolio
%=============================================================================

\documentclass[12pt,letterpaper]{article}
\usepackage[margin=1in]{geometry}
\usepackage{amsmath,amssymb}
\usepackage{enumitem}
\usepackage{hyperref}
\usepackage{tcolorbox}
\usepackage{booktabs}

\title{New Patent Proposals Arising from\\
Patent \#11 (EM Coupling SRF) Research\\[6pt]
\large Three Gap-Filling Applications for the DFD Patent Portfolio}
\author{Gary Thomas Alcock}
\date{\today}

\begin{document}
\maketitle

\section*{Executive Summary}

Research supporting Patent \#11 (Systems and Methods for Electromagnetic
Coupling, Actuation, and Application of Scalar Refractive Fields) has
identified three significant gaps in the existing DFD patent portfolio
(Patents \#1--15). Each gap represents a distinct technology domain not
adequately covered by current filings but essential to the complete
commercialization of EM$\to\psi$ coupling technology. The three proposed
patents are:

\begin{enumerate}
\item \textbf{Patent \#16}: $\psi$-Mode Spectroscopy and Resonance
Identification Systems
\item \textbf{Patent \#17}: Superconductor-Enhanced EM$\to\psi$ Coupling
Devices and Methods
\item \textbf{Patent \#18}: EM$\to\psi$-Mediated Biological and Medical
Sensing Systems
\end{enumerate}

%=============================================================================
\section{Patent Proposal \#16: $\psi$-Mode Spectroscopy and\\
Resonance Identification Systems}
%=============================================================================

\subsection{Gap Analysis}

All DFD actuation methods (driven resonance, parametric amplification,
impulsive excitation) depend critically on knowledge of the $\psi$-mode
natural frequency $\Omega_\psi$. The driven-mode equation:
\[
\ddot{q} + 2\gamma_\psi\dot{q} + \Omega_\psi^2 q
= (\lambda-1)F_{\rm EM}(t)
\]
shows that resonant enhancement occurs only when $2\omega \approx \Omega_\psi$
(driven channel) or $2\omega \approx 2\Omega_\psi$ (parametric channel).
Missing the resonance by even a modest factor dramatically reduces
coupling efficiency.

\textbf{No existing patent covers methods for determining $\Omega_\psi$.}
This is analogous to having radio transmitter/receiver patents without
covering the spectrum analyzer. Without this patent, competitors could
develop $\psi$-mode identification technology independently and control
the essential ``tuning'' step for all DFD devices.

\subsection{Prior Art Distinction}

Standard electromagnetic spectroscopy (RF, microwave, optical) identifies
EM resonance modes. Standard gravitational wave detection (LIGO, pulsar
timing) searches for tensor perturbations. $\psi$-mode spectroscopy is
distinct because:
\begin{itemize}
\item It targets \emph{scalar} refractive field modes, not EM or tensor GW modes
\item It uses EM actuation as the \emph{pump} and non-EM observables
(clock rates, accelerometry, interferometric phase) as the \emph{probe}
\item The spectral signatures are in the EM stress tensor trace
($B^2 - E^2/c^2$) domain, not the EM field amplitude domain
\end{itemize}

\subsection{Field of the Invention}

Systems and methods for identifying, characterizing, and cataloging
natural frequencies ($\Omega_\psi$), damping rates ($\gamma_\psi$),
spatial mode profiles ($u(\mathbf{r})$), and coupling overlaps ($G$, $H$)
of scalar refractive field ($\psi$) modes in laboratory and natural
environments.

\subsection{Proposed Claims}

\begin{tcolorbox}[colback=blue!5!white, colframe=blue!50!black,
title=Patent \#16 -- Proposed Claims]

\textbf{Family A --- Swept-Frequency $\psi$-Spectroscopy}
\begin{enumerate}
\item \textbf{(Independent, Method)} A method for identifying scalar
refractive field ($\psi$) mode resonance frequencies comprising:
\begin{enumerate}[label=(\alph*)]
\item applying electromagnetic energy at a controlled frequency $\omega$
to a region of space using at least one EM source;
\item sweeping the EM drive frequency such that $2\omega$ traverses a
candidate range for $\Omega_\psi$;
\item monitoring response of at least one $\psi$-sensitive detector
selected from: optical cavity frequency, atomic clock transition
frequency, matter-wave interferometer phase, accelerometer output,
or frequency-comb timing link;
\item identifying peaks in the detector response as a function of
$2\omega$ corresponding to $\psi$-mode resonances.
\end{enumerate}

\item The method of claim 1, wherein the EM source employs TE+TM
mode superposition to maximize the driven overlap integral $G$.

\item The method of claim 1, wherein the EM source employs amplitude
modulation at $2\omega$ to access the parametric pumping channel.
\end{enumerate}

\textbf{Family B --- Broadband/Impulsive $\psi$-Spectroscopy}
\begin{enumerate}[resume]
\item A method for mapping the $\psi$-mode spectrum comprising:
\begin{enumerate}[label=(\alph*)]
\item applying a broadband electromagnetic impulse (capacitor discharge,
laser pulse, or electromagnetic pulse) to excite $\psi$ modes across
a wide frequency band;
\item recording the temporal response (ringdown) using at least one
$\psi$-sensitive detector;
\item performing Fourier or wavelet analysis on the ringdown signal to
extract $\Omega_\psi$ and $\gamma_\psi$ for each identified mode.
\end{enumerate}

\item The method of claim 4, further comprising spatial scanning of
the detector array to reconstruct the mode profile $u(\mathbf{r})$
for each identified resonance.
\end{enumerate}

\textbf{Family C --- $\psi$-Spectrum Mapping System}
\begin{enumerate}[resume]
\item \textbf{(Independent, System)} A $\psi$-mode spectroscopy system
comprising:
\begin{enumerate}[label=(\alph*)]
\item at least one tunable or broadband EM source configured for
$\psi$-mode excitation;
\item an array of $\psi$-sensitive detectors;
\item a signal processing unit configured to: correlate detector
responses with EM source parameters, compute spectral decomposition,
identify $\psi$-mode resonances, and generate a $\psi$-mode frequency
map.
\end{enumerate}

\item The system of claim 6 configured for deployment in natural
environments including magnetospheric, solar coronal, or astrophysical
plasma regions.

\item The system of claim 6 further comprising a feedback controller
that automatically tunes EM source parameters to lock onto identified
$\psi$-mode resonances for sustained excitation.
\end{enumerate}

\textbf{Family D --- $\psi$-Mode Database and Prediction}
\begin{enumerate}[resume]
\item A non-transitory computer-readable medium storing:
\begin{enumerate}[label=(\alph*)]
\item a database of measured $\psi$-mode parameters
($\Omega_\psi$, $\gamma_\psi$, $u(\mathbf{r})$, $G$, $H$) indexed
by environment type, geometry, and material composition;
\item algorithms for predicting $\psi$-mode spectra in new environments
based on known EM energy distributions and boundary conditions;
\item optimization routines for computing EM source configurations that
maximize coupling to a target $\psi$ mode.
\end{enumerate}

\item The medium of claim 9, wherein the prediction algorithms
incorporate the area-ratio design law:
\[
|\lambda-1|_{\rm min} = \frac{\pi\gamma_\psi}{c_s U_0 m}
\frac{A_\psi^2}{\kappa_{\rm eff}A_{\rm cav,tot}}
\]
to compute minimum detectable coupling for proposed experimental
configurations.
\end{enumerate}
\end{tcolorbox}

\subsection{Commercial Significance}

This patent controls the essential ``tuning'' technology for all
EM$\to\psi$ devices. Any entity developing DFD-based technology must
identify $\Omega_\psi$ before efficient actuation is possible. This
patent ensures that the spectroscopy/characterization step remains
within the DFD patent family.

%=============================================================================
\section{Patent Proposal \#17: Superconductor-Enhanced\\
EM$\to\psi$ Coupling Devices and Methods}
%=============================================================================

\subsection{Gap Analysis}

Patent \#11 mentions superconductors as EM sources (claim family A,
embodiment 2: LC/cavity), but does not specifically cover the unique
advantages of superconducting systems for EM$\to\psi$ coupling:
\begin{itemize}
\item \textbf{Macroscopic quantum coherence}: All Cooper pairs act
in phase, potentially providing $N^{1/2}$ enhancement
\item \textbf{Flux quantization}: Discrete flux quanta $\Phi_0 = h/(2e)$
provide natural frequency combs when vortices move
\item \textbf{London moment}: $B_L = -2m_e\Omega/e$ provides rotation-frequency-locked
magnetic fields
\item \textbf{Meissner shielding}: Perfect diamagnetism creates extreme
field gradients at surfaces
\item \textbf{Josephson effects}: AC Josephson frequency $\omega_J = 2eV/\hbar$
provides voltage-tunable EM oscillation at arbitrary frequencies
\item \textbf{Extreme $Q$ factors}: Superconducting cavities achieve
$Q > 10^{11}$, maximizing stored energy per watt of input
\end{itemize}

The Tajmar experiments (Sec.~\ref{sec:tajmar} of the main paper) and
the DFD Cooper-pair mass anomaly predictions provide additional
motivation.

\subsection{Field of the Invention}

Devices and methods exploiting superconducting quantum coherence,
flux quantization, Josephson effects, London moment, and extreme
quality factors for enhanced electromagnetic actuation of the scalar
refractive field $\psi$.

\subsection{Proposed Claims}

\begin{tcolorbox}[colback=blue!5!white, colframe=blue!50!black,
title=Patent \#17 -- Proposed Claims]

\textbf{Family A --- Cooper-Pair Coherent Enhancement}
\begin{enumerate}
\item \textbf{(Independent, Method)} A method for enhanced $\psi$
actuation comprising:
\begin{enumerate}[label=(\alph*)]
\item providing a superconducting body containing $N$ Cooper pairs
in a coherent quantum state;
\item applying electromagnetic energy to the superconductor such that
the coherent response of the Cooper-pair condensate amplifies the
EM$\to\psi$ coupling by a factor scaling with $N$ or $\sqrt{N}$;
\item detecting the resulting $\psi$ perturbation using at least one
sensor.
\end{enumerate}

\item The method of claim 1, wherein the superconductor is a type-II
material (YBCO, MgB$_2$, or iron-based) and the EM actuation is
achieved through controlled motion of Abrikosov vortices at
frequencies matched to $\Omega_\psi$.

\item The method of claim 1, wherein the superconductor is a type-I
material (Nb, Pb, Al) operated in the Meissner state for maximum
coherence.
\end{enumerate}

\textbf{Family B --- Josephson Junction Arrays}
\begin{enumerate}[resume]
\item \textbf{(Independent, Device)} A $\psi$-actuation device comprising
a Josephson junction array configured to generate phase-coherent EM
oscillations at a frequency $\omega_J = 2eV/\hbar$ that is tunable
by applied voltage $V$ to match $\Omega_\psi/2$ or $\Omega_\psi$,
thereby maximizing resonant EM$\to\psi$ coupling.

\item The device of claim 4, comprising a series array of $N_J$
junctions providing coherent power enhancement of factor $N_J^2$
at the target frequency.

\item The device of claim 4, further comprising a phase-locked loop
that locks $\omega_J$ to a detected $\psi$-mode frequency using
feedback from a $\psi$-sensitive sensor.
\end{enumerate}

\textbf{Family C --- Rotating Superconductor Systems}
\begin{enumerate}[resume]
\item A method of $\psi$ actuation using a rotating superconductor
comprising:
\begin{enumerate}[label=(\alph*)]
\item rotating a superconducting body at angular velocity $\Omega$
to generate London-moment magnetic field $B_L = -2m_e\Omega/e$;
\item modulating $\Omega$ such that $\dot{B}_L$ produces EM time
variation at frequencies matching $\psi$-mode resonances;
\item detecting $\psi$ perturbations using at least one of:
accelerometer, gyroscope, interferometer, or atomic clock.
\end{enumerate}

\item The method of claim 7, further comprising a second,
counter-rotating superconductor for differential measurement and
common-mode noise rejection.
\end{enumerate}

\textbf{Family D --- SRF Cavity Optimization}
\begin{enumerate}[resume]
\item A superconducting radio-frequency (SRF) cavity system optimized
for $\psi$ actuation, comprising:
\begin{enumerate}[label=(\alph*)]
\item an SRF cavity with $Q > 10^{10}$ operated at a frequency $\omega$
such that $2\omega \approx \Omega_\psi$;
\item a geometry configured with broken symmetry (asymmetric irises,
TE+TM mode superposition, or elliptical cross-section) to maximize
the driven overlap $G$;
\item stored energy $U_0 > 1$~MJ for maximum parametric coupling.
\end{enumerate}

\item The system of claim 9, further comprising an amplitude modulation
system providing controlled modulation depth $m$ at $2\omega$ for
parametric amplification of $\psi$ modes.
\end{enumerate}

\textbf{Family E --- Cryogenic $\psi$-Detection}
\begin{enumerate}[resume]
\item A cryogenic $\psi$-detection system comprising superconducting
sensors (SQUIDs, kinetic inductance detectors, or nanomechanical
resonators) configured to detect $\psi$ perturbations at sensitivities
approaching $\delta\psi \sim 10^{-20}$ through their effect on
local magnetic flux, kinetic inductance, or resonator frequency.
\end{enumerate}
\end{tcolorbox}

\subsection{Commercial Significance}

Superconducting technology represents the highest-performance pathway
to EM$\to\psi$ actuation. Superconducting accelerator cavities
(CERN, DESY, JLab) already store megajoule-scale EM energy at
$Q > 10^{10}$---the exact regime where DFD predicts maximum
$\psi$-coupling sensitivity. This patent ensures that the
superconductor-specific optimization pathway is protected.

%=============================================================================
\section{Patent Proposal \#18: EM$\to\psi$-Mediated\\
Biological and Medical Sensing Systems}
%=============================================================================

\subsection{Gap Analysis}

The existing DFD patent portfolio covers: cloaking/drag (\#1),
resonators/metamaterials (\#2), quantum control (\#3), energy (\#4),
navigation/timing (\#5), sensors/imaging (\#6), energy harvesting (\#7),
communications (\#8), positioning (\#9), computing (\#10), EM coupling
(\#11), and provisionals (\#12--15).

\textbf{No patent covers biological or medical applications.} This is
a significant commercial gap because:
\begin{itemize}
\item Medical imaging is a \$50B+ annual market
\item Biological tissues have diverse, well-characterized EM properties
\item Tissue-specific EM$\to\psi$ coupling signatures could provide a
fundamentally new imaging contrast mechanism
\item The noninvasive nature (EM input, gravitational/refractive output)
avoids ionizing radiation
\end{itemize}

\subsection{Scientific Basis}

Biological tissues have tissue-specific dielectric properties:
\begin{center}
\begin{tabular}{lcc}
\toprule
Tissue & $\epsilon_r$ (1 MHz) & $\sigma$ (S/m) \\
\midrule
Blood & 3000 & 0.7 \\
Muscle & 1600 & 0.5 \\
Fat & 50 & 0.02 \\
Bone & 250 & 0.02 \\
Tumor & 4000+ & 0.8+ \\
\bottomrule
\end{tabular}
\end{center}

When EM energy passes through tissue, the local EM energy density
$U_{\rm EM}(\mathbf{r}) = (\epsilon_r\epsilon_0 E^2 + B^2/\mu_0)/2$
is tissue-dependent. Through EM$\to\psi$ coupling:
\[
\delta\psi(\mathbf{r}) \propto (\lambda-1)\,\epsilon_r(\mathbf{r})\,E^2(\mathbf{r}),
\]
providing a $\psi$-contrast that maps tissue dielectric properties.

The dual-sector bracket further enhances contrast:
\[
\mathcal{B}(\mathbf{r}) = \frac{B^2}{\mu(\psi)} - \epsilon(\psi)E^2
\propto \text{tissue-dependent},
\]
giving different signatures for electric-dominant vs.\ magnetic-dominant
tissue responses.

\subsection{Field of the Invention}

Systems and methods for noninvasive biological and medical sensing,
imaging, and diagnostics using electromagnetic actuation of scalar
refractive field ($\psi$) perturbations, wherein tissue-specific
EM energy distributions produce characteristic $\psi$ signatures
detectable by external sensors.

\subsection{Proposed Claims}

\begin{tcolorbox}[colback=blue!5!white, colframe=blue!50!black,
title=Patent \#18 -- Proposed Claims]

\textbf{Family A --- $\psi$-Contrast Tissue Sensing}
\begin{enumerate}
\item \textbf{(Independent, Method)} A method of noninvasive biological
sensing comprising:
\begin{enumerate}[label=(\alph*)]
\item applying electromagnetic energy at one or more frequencies to a
biological tissue volume;
\item detecting $\psi$ perturbations caused by tissue-specific
electromagnetic energy density distributions using at least one
$\psi$-sensitive sensor external to the tissue;
\item computing tissue characterization parameters from the detected
$\psi$ signature, wherein different tissue types produce
distinguishable $\psi$ coupling profiles due to their different
dielectric properties.
\end{enumerate}

\item The method of claim 1, wherein the EM energy is applied at
multiple frequencies spanning $10^3$--$10^{10}$~Hz to generate a
frequency-dependent $\psi$-response spectrum characteristic of tissue
composition.

\item The method of claim 1, further comprising comparison of the
detected $\psi$ signature against a database of known tissue $\psi$-coupling
profiles to identify tissue type and detect anomalies.
\end{enumerate}

\textbf{Family B --- $\psi$-Tomographic Imaging}
\begin{enumerate}[resume]
\item \textbf{(Independent, System)} A $\psi$-contrast imaging system
comprising:
\begin{enumerate}[label=(\alph*)]
\item an EM source array configured for controlled illumination of a
biological tissue volume from multiple angles;
\item a $\psi$-sensitive detector array surrounding the tissue volume;
\item a computational reconstruction unit configured to generate
tomographic images of tissue EM$\to\psi$ coupling coefficient
distributions.
\end{enumerate}

\item The system of claim 4, wherein the $\psi$-sensitive detectors
comprise precision atomic clocks, optical cavity frequency monitors,
or matter-wave interferometers.

\item The system of claim 4, further comprising dual-sector (TE/TM)
illumination for independent measurement of electric-sector and
magnetic-sector tissue responses.
\end{enumerate}

\textbf{Family C --- Pathology Detection}
\begin{enumerate}[resume]
\item A method for detecting pathological tissue conditions comprising:
\begin{enumerate}[label=(\alph*)]
\item performing $\psi$-contrast sensing according to claim 1 on a
tissue region of interest;
\item comparing the $\psi$ signature against baseline measurements of
healthy tissue;
\item identifying deviations exceeding a threshold as indicative of
pathological changes in tissue dielectric properties.
\end{enumerate}

\item The method of claim 7, wherein the pathological conditions
detected include at least one of: tumor tissue (elevated dielectric
constant), ischemic tissue (altered conductivity), calcification
(reduced dielectric constant), or edema (elevated water content).
\end{enumerate}

\textbf{Family D --- Wearable/Continuous Monitoring}
\begin{enumerate}[resume]
\item A wearable biological monitoring device comprising:
\begin{enumerate}[label=(\alph*)]
\item a miniaturized EM source configured for safe, low-power
continuous tissue illumination;
\item at least one miniaturized $\psi$-sensitive sensor;
\item a processor configured for continuous computation of
tissue $\psi$-coupling parameters and alerting on significant
deviations.
\end{enumerate}

\item The device of claim 9, configured for continuous blood glucose
monitoring through correlation of blood dielectric changes with
$\psi$-coupling signature variations.
\end{enumerate}
\end{tcolorbox}

\subsection{Commercial Significance}

Medical/biological sensing based on $\psi$-contrast represents a
potentially transformative imaging modality:
\begin{itemize}
\item No ionizing radiation (unlike X-ray/CT)
\item No strong magnetic fields required (unlike MRI)
\item Tissue contrast from fundamentally different physical mechanism
\item Potential for real-time, wearable monitoring
\item Enormous market (\$50B+ medical imaging, \$30B+ diagnostics)
\end{itemize}

This patent secures the medical/biological domain for the DFD
technology family before competitors can identify and claim this space.

%=============================================================================
\section{Filing Priority and Strategy}
%=============================================================================

\begin{table}[h]
\centering
\caption{Recommended filing priority for new patents.}
\begin{tabular}{clcc}
\toprule
Priority & Patent & Market Size & Defensive Value \\
\midrule
1 & \#16 ($\psi$-Spectroscopy) & Medium & \textbf{Critical} \\
2 & \#17 (Superconductor) & Large & High \\
3 & \#18 (Medical/Bio) & Very Large & Medium \\
\bottomrule
\end{tabular}
\end{table}

\textbf{Rationale:}
\begin{itemize}
\item Patent \#16 should be filed first because it is the enabling
technology for \emph{all} other DFD patents. Without $\psi$-mode
identification, no device can be optimally tuned.
\item Patent \#17 should follow because superconducting systems
represent the highest-performance implementation pathway and the
most likely route to first detection of $|\lambda-1| \neq 0$.
\item Patent \#18 can be filed later as it depends on successful
demonstration of $\psi$ detection (i.e., after Patents \#16--17
yield results), but should be filed before any public disclosure
of biological applications.
\end{itemize}

All three patents should cross-reference Patent \#11 as the parent
umbrella and claim priority from the August 16, 2025 provisional
filings where applicable.

\end{document}
